%%
%% Copyright 2019-2020 Elsevier Ltd
%%
%% This file is part of the 'CAS Bundle'.
%% --------------------------------------
%%
%% It may be distributed under the conditions of the LaTeX Project Public
%% License, either version 1.2 of this license or (at your option) any
%% later version.  The latest version of this license is in
%%    http://www.latex-project.org/lppl.txt
%% and version 1.2 or later is part of all distributions of LaTeX
%% version 1999/12/01 or later.
%%
%% The list of all files belonging to the 'CAS Bundle' is
%% given in the file `manifest.txt'.
%%
%% Template article for cas-sc documentclass for
%% double column output.

\documentclass[a4paper,fleqn]{cas-sc}
\usepackage{njfutitlepage}
\usepackage[authoryear]{natbib}
\usepackage{csquotes}
\usepackage[UTF8]{ctex}
\usepackage{gbt7714}
\usepackage{makecell}
\usepackage{graphicx}
\usepackage{booktabs}
\usepackage{multirow}
\usepackage{url}
\usepackage{listings}
\usepackage{xcolor}
\usepackage{hyperref}
\usepackage{amsmath}
\usepackage{amssymb}
\usepackage{placeins}

\newcommand{\refSec}[1]{\hyperref[#1]{Section~\ref*{#1}}}
\newcommand{\refFig}[1]{\hyperref[#1]{Fig.~\ref*{#1}}}
\newcommand{\refApp}[1]{\hyperref[#1]{Appendix~\ref*{#1}}}
\renewcommand{\figurename}{Fig.}


\lstset{
    language=bash,
    basicstyle=\footnotesize\ttfamily,
    keywordstyle=\color{blue}\bfseries,
    commentstyle=\color{gray},
    stringstyle=\color{red},
    numbers=left,
    numberstyle=\tiny\color{gray},
    stepnumber=1,
    numbersep=8pt,
    backgroundcolor=\color{gray!10},
    frame=single,
    tabsize=2,
    captionpos=b,
    breaklines=true,
    breakatwhitespace=false,
    escapeinside={\%*}{*)}
}

%%%Author definitions
\def\tsc#1{\csdef{#1}{\textsc{\lowercase{#1}}\xspace}}
\tsc{WGM}
\tsc{QE}
\tsc{EP}
\tsc{PMS}
\tsc{BEC}
\tsc{DE}
%%%

\begin{document}
\let\WriteBookmarks\relax
\def\floatpagepagefraction{1}
\def\textpagefraction{.001}

\shortauthors{7组}
\shorttitle{}
\title[mode = title]{FeHALS 航路规划与激光仿真系统 API 技术文档}

\Course{\enquote{智能+}GIS设计与开发}
\Major{测绘工程}
\Grade{2023级}
\Group{7 组}

\Member[2310604116]{王 孜悠}
\Member[2310604109]{梅 康凯}
\Member[2310604115]{唐 如意}
\Member[2310604122]{杨 杰瑞}
\Member[2310604117]{温 璞}
\Member[2110604112]{刘 继明}
\Member[2110604106]{李 英杰}
\Member[2310604103]{杜 仲宇}
\Member[2310604113]{施 祺}
\Member[2310604118]{向 宇鸣}

\Teacher{隋 铭明}
\Teacher{那 嘉明}
\Teacher{陈 动}
\Teacher{朱 杰}


\MakeNJFUTitlePage[UseDocTitle]
\tableofcontents
\newpage

\section{系统概述}

\subsection{项目背景}
FeHALS（Frontend of HELIOS++ for Airborne Laser Scanning）是一个基于Web的三维可视化航路规划与激光仿真系统，旨在解决HELIOS++原生命令行交互复杂、缺乏三维可视化航路规划能力的问题。本系统为航测工程师、遥感研究人员以及相关专业师生提供一个直观易用的平台，支持从场景构建、航路设计、仿真执行到结果可视化的全流程操作。

\subsection{系统架构}
系统采用浏览器/服务器（B/S）三层架构：
\begin{itemize}
    \item \textbf{前端层}：基于Vue 3与Three.js构建，负责三维场景渲染、航点交互编辑、参数配置与结果展示
    \item \textbf{后端层}：基于FastAPI实现，负责API服务、文件生成、仿真调度与数据处理
    \item \textbf{仿真引擎层}：HELIOS++激光扫描仿真引擎，独立运行并由后端调用
\end{itemize}

\begin{figure}[pos=!htbp]
    \centering
    \caption{FeHALS系统架构与数据流}
    \includegraphics[width=0.8\textwidth]{figs/gernal_fehals_arch.pdf}
    \label{fig:architecture}
\end{figure}

\section{相关技术}\label{sec:tech}

\subsection{HELIOS++仿真框架}
HELIOS++是一款基于C++的开源激光雷达仿真框架，采用光线追踪技术对激光脉冲在三维场景中的传播过程进行建模，能够模拟机载、地面、车载与无人机等多种平台，支持全波形（full-waveform）输出与LAS、LAZ、XYZ等多种点云格式。其仿真配置以XML文件组织：场景（scene）文件定义参与扫描的三维地物模型；平台（platform）与扫描器（scanner）目录文件定义飞行平台与扫描仪的物理参数；扫描任务（survey）文件将场景、平台、扫描器与航迹（trajectory）关联起来，并设定脉冲频率、扫描频率与扫描角等核心参数。

\subsection{Web三维可视化与Three.js}
Three.js是一个基于WebGL \citep{webgl} 的开源JavaScript三维图形库，封装了场景（Scene）、相机（Camera）、渲染器（Renderer）与网格（Mesh）等核心对象，支持OBJ、GLTF、STL等多种三维模型格式的加载，并提供了轨道控制（OrbitControls）、拖拽控制（DragControls）与射线拾取（Raycaster）等交互组件。借助Three.js，可在浏览器中构建轻量级的三维可视化场景，实现模型展示、点云渲染与鼠标交互拾取，为航路规划提供直观的操作界面。

\subsection{Vue 3前端框架}
Vue 3是一个渐进式JavaScript框架，采用Composition API提供更好的代码组织和类型支持。其核心特性包括响应式系统、组件化开发和虚拟DOM。本系统前端基于Vue 3构建，使用Pinia进行状态管理，通过组件化方式组织用户界面，实现了高效的3D交互和参数配置功能。

\subsection{FastAPI与后端异步任务}
FastAPI是一个基于Python的高性能Web框架，基于类型注解与Pydantic \citep{pydantic} 实现自动的数据校验与接口文档生成，并原生支持异步编程与WebSocket。本系统后端基于FastAPI构建REST API \citep{fielding2000rest} 与WebSocket日志通道，通过asyncio子进程异步调用HELIOS++可执行文件，实时捕获其标准输出并以WebSocket推送到前端，实现仿真过程的实时监控。

\section{系统需求分析与设计}\label{sec:design}

\subsection{用户需求分析}
本系统的目标用户为航测工程师、遥感研究人员以及相关专业师生。其核心需求可归纳为三个方面：一是直观的航路规划，用户希望在不编写配置文件的情况下，通过三维场景中的鼠标交互即可完成航点设定与航线编辑；二是快速的仿真执行，用户能够通过图形化界面配置平台与扫描仪参数，一键调用HELIOS++引擎并实时查看运行状态；三是便捷的结果可视化，仿真生成的点云能够自动加载到三维场景中，并支持按高度或强度着色以及尺寸、透明度的调节。

\subsection{系统总体架构}
系统采用浏览器/服务器（B/S）架构，分为前端、后端与仿真引擎三层，其总体架构如图\ref{fig:architecture}所示。前端基于Vue 3构建，负责三维场景渲染、航点交互编辑、参数配置与结果展示，通过HTTP \citep{rfc9110,rfc9113,rfc9114} 与WebSocket与后端通信；后端基于FastAPI，负责接收前端请求，生成HELIOS++所需的航迹文件与XML配置文件，调用HELIOS++可执行文件执行仿真，并解析与返回点云结果；仿真引擎层为HELIOS++，独立于前后端，由后端以子进程方式调用。各模块间通过明确的REST API与WebSocket消息格式进行数据交换，形成完整的工作流。

\subsection{功能优先级划分}
结合开发环境与工程约束，系统功能划分为三个优先级：核心功能（MVP）包括三维场景渲染、模型加载、交互式航点编辑、航迹导出、仿真参数配置、仿真执行、日志控制台与点云渲染；重要功能包括场景要素属性编辑、自动航线生成、航高自动计算、仿真预设方案、点云特征统计与点云下载；拓展功能包括仿真进度实时反馈、点云覆盖度分析、多任务队列管理与仿真过程动画。

\section{系统设计与实现}\label{sec:impl}

\subsection{系统模块划分}
系统实现按前后端分层组织，前端基于Vue 3与Node.js \citep{nodejs} 生态构建，目录下包含组件（components）、状态管理（stores）与组合式函数（composables）三部分：组件层包含三维场景Scene3D.vue、参数配置面板ControlPanel.vue、航点列表WaypointList.vue与日志控制台LogConsole.vue；stores层基于Pinia维护场景（scene）、航点（waypoints）与仿真（simulation）三份响应式状态；composables层封装Three.js场景管理（useThreeScene）、航点交互（useWaypoints）与后端接口调用（useHeliosAPI）。

后端基于FastAPI构建，api目录下包含REST路由routes.py与WebSocket端点websocket.py，models目录定义Pydantic数据模型schemas.py，services目录封装航迹生成trajectory\_generator.py、配置生成config\_generator.py、HELIOS++调用helios\_service.py与点云解析pointcloud\_parser.py。

\subsection{三维场景渲染与航点交互}
三维场景基于Three.js构建，Scene3D.vue负责初始化场景、相机、渲染器与灯光，并添加地面网格与拾取平面。场景使用Z-up坐标系（与HELIOS++一致，z为高程轴），通过OrbitControls实现相机的旋转、平移与缩放；模型加载采用OBJLoader、GLTFLoader与STLLoader，加载完成后根据模型包围盒自动调整相机视角，并基于模型顶点包围盒推断其up轴（Y-up模型自动旋转至Z-up）。航点交互通过自绘拖拽实现：鼠标左键点击经Raycaster与地面平面求交，在交点处添加球体表示的航点（恒为z=0贴地），相邻航点之间以白色折线连接；拖拽航点时锁定相机（`controls.enabled=false`），确保不干扰视角操作；通过Delete键或右键菜单删除航点。

\subsection{前端组件与状态管理}
侧边栏采用多Tab布局，包含仿真参数、点云、模型列表、航迹与设置五个Tab。ControlPanel.vue提供仿真参数表单，分为「载体参数」与「传感器参数」两个模块，各参数旁显示有效范围，只读参数以灰色禁用显示；参数数据来源于从HELIOS++扫描器与平台目录文件中提取的规格库，切换平台类型时自动重置参数至默认值。WaypointList.vue以列表形式展示航点序号与坐标，支持单个删除与整体清空。LogConsole.vue展示分级日志（INFO/WARNING/ERROR），支持自动滚动与清空。

三个Pinia store分别维护场景模型与点云渲染参数（scene）、航点坐标数组（waypoints）以及仿真任务状态、进度、日志与结果（simulation）。前端通过useHeliosAPI.js中的Axios实例调用后端REST接口，并通过connectLogWS建立WebSocket连接以接收仿真日志。

\subsection{后端API接口设计}
后端基于FastAPI实现，提供完整的RESTful API接口，主要分为以下六类：

\begin{enumerate}
    \item \textbf{模型管理}：上传/列表/删除三维模型，支持OBJ、GLTF、GLB、STL格式
    \item \textbf{航迹管理}：从航点生成HELIOS++原生.trj航迹文件，支持弓字形自动生成
    \item \textbf{配置管理}：生成HELIOS++ survey XML与scene XML配置文件
    \item \textbf{仿真执行}：异步调用HELIOS++子进程，实时推送日志，支持状态查询与取消
    \item \textbf{结果管理}：返回点云统计（点数/边界/高程/强度/直方图）、渲染抽样与文件下载
    \item \textbf{环境诊断}：检测HELIOS++可执行文件、资源目录、静态工作目录的就绪状态
\end{enumerate}

\begin{table}[pos=!htbp]
    \centering
    \caption{FeHALS REST API接口列表}
    \label{tab:api-interfaces}
    \begin{tabular}{@{} lll @{}}
        \toprule
        \textbf{接口组} & \textbf{HTTP方法}                          & \textbf{功能说明} \\
        \midrule
        模型管理         & GET /api/models                          & 获取所有上传的模型列表   \\
                     & POST /api/models/upload                  & 上传3D模型文件      \\
                     & DELETE /api/models/\{id\}                & 删除指定模型        \\
        航迹管理         & POST /api/trajectory/generate            & 从航点生成航迹文件     \\
                     & GET /api/trajectories                    & 获取所有航迹列表      \\
                     & GET /api/trajectories/download/\{id\}    & 下载航迹文件        \\
        配置管理         & POST /api/config/generate                & 生成仿真配置文件      \\
        仿真执行         & POST /api/simulation/run                 & 执行HELIOS++仿真  \\
                     & GET /api/simulation/status/\{task\_id\}  & 获取仿真状态        \\
                     & GET /api/simulation/logs/\{task\_id\}    & 获取仿真日志        \\
                     & POST /api/simulation/cancel/\{task\_id\} & 取消仿真任务        \\
        结果管理         & GET /api/results/\{task\_id\}            & 获取点云数据与统计     \\
                     & GET /api/results/\{task\_id\}/download   & 下载原始结果        \\
        环境诊断         & GET /api/env/diagnose                    & 检查系统环境        \\
        缓存管理         & GET /api/cache                           & 查看各缓存目录大小     \\
                     & DELETE /api/cache/\{type\}               & 分类清理缓存        \\
        \bottomrule
    \end{tabular}
\end{table}

\subsection{航迹与配置文件生成}
trajectory\_generator.py将航点数组转换为HELIOS++原生航迹文件（.trj，CSV格式，列顺序为时间t、roll、pitch、yaw、x、y、z），所有航点统一为恒定飞行高度（仅用x、y定义水平路径）。在轨迹生成过程中，系统按以下步骤计算航迹点：

\begin{enumerate}
    \item \textbf{偏航角计算}：从当前航点至下一航点的方向向量经 $\text{yaw} = \arctan(\Delta x / \Delta y)$ 转换为ARINC 705规范的偏航角，yaw=0°表示正北方向（+Y），yaw=90°表示正东方向（+X），确保扫描线始终正交于航线方向。
    \item \textbf{时间步长计算}：根据目标飞行速度（默认10 m/s）和航段距离计算各航段的时间步长 $dt = \max(1.0, dist / v_{\text{target}})$，时间沿航线累积，使各航段的速度连续且物理合理。
    \item \textbf{转角瞬时转向}：在航点转角处，系统生成同一时间戳、同一位置、不同偏航角的两行记录，第一行保持上一段的偏航角，第二行切换至下一段的偏航角，实现瞬时转向，避免偏航角在航段间渐变造成扫描线偏移。
\end{enumerate}

系统还支持弓字形自动航迹生成：用户在三维场景中点击两个角点定义矩形区域，系统根据设定航线间距自动生成往返覆盖的全覆盖航线，相邻扫描线的X端点次序交替，形成连续往返路径。扫描条带宽度由 $W = 2 \cdot h \cdot \tan(\theta / 2)$ 计算，扫描线间距由 $S = v / f_{\text{scan}}$ 计算，其中 $h$ 为飞行高度，$\theta$ 为总扫描角，$v$ 为飞行速度，$f_{\text{scan}}$ 为扫描频率。

config\_generator.py负责生成HELIOS++所需的两个XML文件：场景文件通过objloader滤镜引用三维模型，并指定模型的up轴（Y-up自动旋转至Z-up）；扫描任务（survey）文件引用场景、平台与扫描器目录，并将平台类型映射为对应的平台目录项（UAV与Airborne均使用`copter\_linearpath`平台、`riegl\_vux-1uav`扫描器），将脉冲频率（kHz）与扫描角（±半角）换算为HELIOS++所需的脉冲频率（Hz）与总扫描角。

\subsection{HELIOS++调用与点云渲染}
helios\_service.py维护仿真任务注册表，通过asyncio.create\_subprocess\_exec以子进程方式调用helios++可执行文件，指定资源目录（--assets）、输出目录（--output）与输出格式参数，实时捕获其标准输出/错误流并解析进度百分比，经WebSocket推送至前端，同时支持超时终止与任务取消。仿真启动前，前端对全部参数进行范围校验（数据来源于从HELIOS++扫描器目录文件中提取的参数规格库），超出有效范围时拒绝执行并给出提示。系统还提供HELIOS++运行环境诊断功能，通过GET /api/env/diagnose接口检测可执行文件存在性、资源目录完整性和静态工作目录就绪状态，在前端设置面板中集中展示。

仿真结束后，pointcloud\_parser.py利用laspy \citep{laspy} 解析LAS/LAZ点云或按文本方式解析XYZ点云，提取坐标与强度、计算空间范围。点云统计信息包括总点数、三维边界、高程均值、标准差、最小/最大、中位数、P5/P95分位数，以及40组高程分布直方图。强度字段存在时同样计算其均值、标准差和范围。对于超过百万级的点云，使用覆盖全文件序列的等间隔索引抽样进行渲染，全量统计基于完整点集计算，确保统计准确性。前端将返回的点云坐标构建为BufferGeometry与Points对象，支持按高度（蓝色至红色渐变）、按强度或固定颜色着色，并可实时调节点尺寸与透明度。

\subsection{项目组织结构}
\begin{figure}[pos=!htbp]
    \centering
    \caption{FeHALS项目目录结构}
    \begin{lstlisting}[language=bash]
FeHALS/
+-- backend/                 # 后端代码
|   +-- app/
|   |   +-- main.py         # FastAPI应用入口
|   |   +-- config.py       # 配置管理
|   |   +-- api/            # API路由
|   |   +-- models/         # 数据模型
|   |   +-- services/       # 业务服务
|   |   +-- static/         # 静态文件目录
|   +-- requirements.txt    # Python依赖
|   +-- run.py             # 后端启动脚本
+-- frontend/              # 前端代码
|   +-- src/
|   |   +-- components/    # Vue组件
|   |   +-- composables/  # 组合式函数
|   |   +-- stores/       # Pinia状态管理
|   |   +-- assets/       # 静态资源
|   +-- package.json
|   +-- vite.config.js
+-- doc/                   # 文档目录
|   +-- src/
|   |   +-- DESCRIPTION.tex
|   |   +-- DESIGNBOOK.tex
|   |   +-- DOCUMENT.tex
|   +-- Makefile           # 构建脚本
+-- run.sh                 # Unix启动脚本
  \end{lstlisting}
    \label{fig:project-structure}
\end{figure}

\FloatBarrier

\section{部署与运行指南}

\subsection{环境要求}
\begin{table}[pos=!htbp]
    \centering
    \caption{系统环境要求}
    \begin{tabular}{@{} ll @{}}
        \toprule
        \textbf{组件} & \textbf{要求}      \\
        \midrule
        Node.js     & >= 16.0.0        \\
        Python      & >= 3.8.0         \\
        npm         & >= 7.0.0         \\
        LaTeX       & XeLaTeX (用于文档编译) \\
        \bottomrule
    \end{tabular}
\end{table}

\subsection{安装步骤}
\begin{enumerate}
    \item \textbf{后端安装}
          \begin{lstlisting}[language=bash]
  # 1. 进入后端目录
  cd backend

  # 2. 创建虚拟环境
  python -m venv venv

  # 3. 激活虚拟环境
  # Linux/Mac
  source venv/bin/activate
  # Windows
  venv\Scripts\activate

  # 4. 安装依赖
  pip install -r requirements.txt

  # 5. 设置环境变量
  export HELIOS_PATH=/path/to/helios++
  export HELIOS_REPO=/path/to/helios++/repo
  export HELIOS_ASSETS=/path/to/assets
  \end{lstlisting}

    \item \textbf{前端安装}
          \begin{lstlisting}[language=bash]
  # 1. 进入前端目录
  cd frontend

  # 2. 安装依赖
  npm install

  # 3. 启动开发服务器
  npm run dev
  \end{lstlisting}
\end{enumerate}

\subsection{启动脚本}
项目提供了启动脚本用于快速启动服务：
\begin{lstlisting}[language=bash,caption=run.sh脚本]
#!/bin/bash

# 停止所有服务
stop_services() {
    echo "停止所有服务..."
    pkill -f "uvicorn" || true
    pkill -f "vite" || true
}

# 启动后端
start_backend() {
    echo "启动后端服务..."
    cd backend
    source venv/bin/activate
    python -m uvicorn app.main:app --reload --host 0.0.0.0 --port 8000 &
    cd ..
}

# 启动前端
start_frontend() {
    echo "启动前端服务..."
    cd frontend
    npm run dev &
    cd ..
}

# 主程序
case "$1" in
    start)
        start_backend
        start_frontend
        echo "服务已启动"
        echo "前端: http://localhost:3000"
        echo "后端: http://localhost:8000"
        echo "API文档: http://localhost:8000/docs"
        ;;
    stop)
        stop_services
        echo "服务已停止"
        ;;
    *)
        echo "用法: $0 {start|stop}"
        exit 1
        ;;
esac
\end{lstlisting}

\subsection{配置说明}
\begin{itemize}
    \item \textbf{后端配置}
          \begin{lstlisting}[language=bash]
  # 环境变量示例
  export FEHALS_HOST=0.0.0.0
  export FEHALS_PORT=8000
  export HELIOS_PATH=/usr/local/bin/helios++
  export HELIOS_REPO=/opt/helios++
  export HELIOS_ASSETS=/opt/helios++/assets
  export FEHALS_SIM_TIMEOUT=300
  export FEHALS_CORS_ORIGINS=http://localhost:3000,http://127.0.0.1:3000
  \end{lstlisting}

    \item \textbf{前端配置}
          \begin{lstlisting}[language=bash]
  // vite.config.js
  export default defineConfig({
    server: {
      proxy: {
        '/api': {
          target: 'http://localhost:8000',
          changeOrigin: true
        },
        '/ws': {
          target: 'ws://localhost:8000',
          ws: true
        }
      }
    }
  })
  \end{lstlisting}
\end{itemize}

\section{API接口详解}\label{sec:api}

\subsection{REST API接口}

\subsubsection{模型管理接口}
模型管理接口支持三维模型的增删查操作。请求与响应格式如下：

\begin{table}[pos=!htbp]
    \centering
    \caption{模型管理API接口}
    \label{tab:api-models}
    \begin{tabular}{@{} llll @{}}
        \toprule
        \textbf{接口路径}               & \textbf{方法} & \textbf{请求体}    & \textbf{响应}                      \\
        \midrule
        \texttt{/api/models}        & GET         & 无               & 模型列表（id/name/url/up）             \\
        \texttt{/api/models/upload} & POST        & multipart: file & \{model\_id, filename, url, up\} \\
        \texttt{/api/models/\{id\}} & DELETE      & 路径参数\{id\}      & \{success: true\}                \\
        \bottomrule
    \end{tabular}
\end{table}

\paragraph{上传模型接口示例}
\begin{lstlisting}[language=bash,caption=上传3D模型]
# 请求
curl -X POST http://localhost:8000/api/models/upload \
  -F "file=@building.obj"

# 响应
{"model_id":"mod_1700000000_abc123","filename":"building.obj",
 "url":"/static/models/mod_1700000000_abc123.obj","up":"z"}
\end{lstlisting}

\paragraph{获取模型列表}
\begin{lstlisting}[language=bash,caption=获取模型列表]
# 请求
curl http://localhost:8000/api/models

# 响应
{"models":[
  {"id":"mod_...","name":"building.obj","url":"/static/models/...","up":"z"}
]}
\end{lstlisting}

\subsubsection{航迹管理接口}
航迹管理接口从航点生成HELIOS++原生.trj文件，支持下载。

\begin{table}[pos=!htbp]
    \centering
    \caption{航迹管理API接口}
    \label{tab:api-trajectory}
    \begin{tabular}{@{} llll @{}}
        \toprule
        \textbf{接口路径}                              & \textbf{方法} & \textbf{请求体}                           & \textbf{响应}                               \\
        \midrule
        \texttt{/api/trajectory/generate}          & POST        & \{waypoints: [[x,y,z],...], altitude\} & \{file\_id, download\_url, point\_count\} \\
        \texttt{/api/trajectories}                 & GET         & 无                                      & \{trajectories: [{id, created}]\}         \\
        \texttt{/api/trajectories/download/\{id\}} & GET         & 路径参数\{id\}                             & .trj文件（文本/纯文本）                            \\
        \bottomrule
    \end{tabular}
\end{table}

\paragraph{生成航迹接口示例}
\begin{lstlisting}[language=bash,caption=生成航迹]
# 请求：三个航点构成L形航线
curl -X POST http://localhost:8000/api/trajectory/generate \
  -H "Content-Type: application/json" \
  -d '{"waypoints":[[0,0,0],[100,0,0],[100,100,0]],"altitude":100}'

# 响应
{"file_id":"traj_...","download_url":"/api/trajectories/download/traj_...",
 "point_count":3}

# 生成的.trj文件内容（CSV）
# 0.00,0,0,0.00,0.000000,0.000000,100.000000
# 10.00,0,0,0.00,100.000000,0.000000,100.000000
# 10.00,0,0,90.00,100.000000,0.000000,100.000000
# 20.00,0,0,90.00,100.000000,100.000000,100.000000
\end{lstlisting}

\subsubsection{仿真执行接口}
仿真执行接口异步调用HELIOS++引擎，任务状态通过轮询或WebSocket获取。

\begin{table}[pos=!htbp]
    \centering
    \caption{仿真执行API接口}
    \label{tab:api-simulation}
    \begin{tabular}{@{} llll @{}}
        \toprule
        \textbf{接口路径}                                & \textbf{方法} & \textbf{请求体}                                     & \textbf{响应}                                 \\
        \midrule
        \texttt{/api/simulation/run}                 & POST        & \{trajectory\_id, config\_id, scene\_model\_id\} & \{task\_id\}                                \\
        \texttt{/api/simulation/status/\{task\_id\}} & GET         & 路径参数\{task\_id\}                                 & \{status, progress, message, result\_file\} \\
        \texttt{/api/simulation/logs/\{task\_id\}}   & GET         & 路径参数\{task\_id\}                                 & \{logs: [string]\}                          \\
        \texttt{/api/simulation/cancel/\{task\_id\}} & POST        & 路径参数\{task\_id\}                                 & \{success: bool\}                           \\
        \bottomrule
    \end{tabular}
\end{table}

\paragraph{启动仿真接口示例}
\begin{lstlisting}[language=bash,caption=启动仿真]
# 1. 先上传模型、生成航迹、生成配置，获得各ID
# 2. 提交仿真任务
curl -X POST http://localhost:8000/api/simulation/run \
  -H "Content-Type: application/json" \
  -d '{"trajectory_id":"traj_...","config_id":"cfg_...",
       "scene_model_id":"mod_..."}'

# 响应
{"task_id":"sim_..."}

# 查询状态
curl http://localhost:8000/api/simulation/status/sim_...
# {"task_id":"sim_...","status":"running","progress":45,"message":"","result_file":null}
\end{lstlisting}

\subsubsection{结果管理接口}
仿真完成后，结果接口返回点云统计与渲染抽样数据。

\begin{table}[pos=!htbp]
    \centering
    \caption{结果管理API接口}
    \label{tab:api-results}
    \begin{tabular}{@{} llll @{}}
        \toprule
        \textbf{接口路径}                               & \textbf{方法} & \textbf{参数}                   & \textbf{响应}                                                    \\
        \midrule
        \texttt{/api/results/\{task\_id\}}          & GET         & 路径参数\{task\_id\}              & \{file\_path, point\_count, bounds, stats, points, intensity\} \\
        \texttt{/api/results/\{task\_id\}/download} & GET         & 路径参数\{task\_id\}, ?format=las & 原始结果文件                                                         \\
        \bottomrule
    \end{tabular}
\end{table}

\paragraph{获取结果接口示例}
\begin{lstlisting}[language=bash,caption=获取结果]
curl http://localhost:8000/api/results/sim_...

# 响应
{"file_path":"/static/results/sim_.../output.xyz",
 "point_count":125000,
 "bounds":{"minx":-10.5,"miny":-15.2,"minz":85.3,
           "maxx":110.8,"maxy":115.6,"maxz":102.1},
 "stats":{"mean_z":93.7,"std_z":4.2,"min_z":85.3,"max_z":102.1,
          "median_z":94.1,"p5_z":88.2,"p95_z":100.5,
          "histogram":[[85,100],[86,250],...]},
 "points":{"x":[...],"y":[...],"z":[...],"intensity":[...]}}
\end{lstlisting}

\subsubsection{环境诊断与缓存管理}

\begin{table}[pos=!htbp]
    \centering
    \caption{环境诊断与缓存管理API}
    \label{tab:api-env-cache}
    \begin{tabular}{@{} llll @{}}
        \toprule
        \textbf{接口路径}                & \textbf{方法} & \textbf{参数}  & \textbf{响应}                                      \\
        \midrule
        \texttt{/api/env/diagnose}   & GET         & 无            & \{helios\_status, assets\_status, dirs\_status\} \\
        \texttt{/api/cache}          & GET         & 无            & \{cache: \{models: \{count, size\}, ...\}\}      \\
        \texttt{/api/cache/\{type\}} & DELETE      & 路径参数\{type\} & \{success, removed\}                             \\
        \bottomrule
    \end{tabular}
\end{table}

\subsection{WebSocket接口}

\subsubsection{日志推送}
仿真任务运行期间，后端通过WebSocket向 \texttt{/ws/logs/\{task\_id\}} 推送实时日志。

\begin{table}[pos=!htbp]
    \centering
    \caption{WebSocket日志消息格式}
    \label{tab:ws-logs}
    \begin{tabular}{@{} lll @{}}
        \toprule
        \textbf{消息类型}   & \textbf{数据字段}         & \textbf{说明}                              \\
        \midrule
        \texttt{log}    & \texttt{timestamp}    & 日志时间戳                                    \\
                        & \texttt{level}        & 日志级别（INFO/WARNING/ERROR）                 \\
                        & \texttt{message}      & 日志消息内容                                   \\
                        & \texttt{progress}     & 进度百分比（0-100）                             \\
        \midrule
        \texttt{status} & \texttt{task\_id}     & 任务ID                                     \\
                        & \texttt{status}       & 任务状态（running/completed/failed/cancelled） \\
                        & \texttt{progress}     & 当前进度                                     \\
                        & \texttt{error}        & 错误信息（如果有）                                \\
        \midrule
        \texttt{result} & \texttt{task\_id}     & 任务ID                                     \\
                        & \texttt{result\_path} & 结果文件路径                                   \\
                        & \texttt{point\_count} & 点云总数                                     \\
                        & \texttt{file\_size}   & 结果文件大小                                   \\
        \bottomrule
    \end{tabular}
\end{table}

\paragraph{WebSocket连接示例}
\begin{lstlisting}[language=Java,caption=WebSocket客户端连接]
// 前端 useHeliosAPI.js 中的 WebSocket 连接逻辑
const ws = new WebSocket(`ws://${host}/ws/logs/${taskId}`);
ws.onmessage = (evt) => {
  const msg = JSON.parse(evt.data);
  switch (msg.type) {
    case 'log':    console.log(`[${msg.level}] ${msg.message}`); break;
    case 'status': updateStatus(msg.status, msg.progress); break;
    case 'result': loadResult(msg.task_id); break;
  }
};
\end{lstlisting}

\subsection{数据模型定义}

后端的请求与响应数据结构由 Pydantic 模型定义，源文件位于 \texttt{backend/app/models/schemas.py}：

\lstinputlisting[language=python, caption={backend/app/models/schemas.py}]{../backend/app/models/schemas.py}

\section{前端实现详解}\label{sec:frontend}

\subsection{Vue 3应用架构}

\subsubsection{组件结构}
前端采用Vue 3 Composition API构建，主要组件包括：
\begin{lstlisting}[language=bash,caption=前端组件结构]
frontend/src/
+-- components/          # Vue组件
|   +-- Scene3D.vue     # 3D场景渲染组件
|   +-- ControlPanel.vue  # 参数配置面板
|   +-- WaypointList.vue  # 航点列表组件
|   +-- LogConsole.vue    # 日志控制台
|   +-- ModelList.vue     # 模型管理组件
|   +-- PointCloudPanel.vue # 点云渲染面板
|   +-- SettingsPanel.vue # 设置面板
+-- stores/            # Pinia状态管理
|   +-- scene.js       # 场景状态
|   +-- waypoints.js   # 航点状态
|   +-- simulation.js  # 仿真状态
+-- composables/      # 组合式函数
    +-- useThreeScene.js  # Three.js场景管理
    +-- useWaypoints.js   # 航点交互逻辑
    +-- useHeliosAPI.js   # API调用封装
    +-- useBowtie.js      # 弓字形航线生成
\end{lstlisting}

\subsubsection{3D场景组件实现}
以下为 \texttt{Scene3D.vue} 的核心实现，包含场景初始化、模型加载与航点渲染的调度逻辑：

\lstinputlisting[language=Java, caption={frontend/src/components/Scene3D.vue}]{../frontend/src/components/Scene3D.vue}

\subsection{Three.js集成与3D交互}

\subsubsection{场景管理组合式函数}
以下为 \texttt{useThreeScene.js} 的核心实现，封装了Three.js场景、相机、渲染器、模型加载与航点交互逻辑：

\lstinputlisting[language=Java, caption={frontend/src/composables/useThreeScene.js}]{../frontend/src/composables/useThreeScene.js}

\subsubsection{航点交互系统}
以下为 \texttt{useWaypoints.js} 的实现，作为面向组件的航点交互门面：

\lstinputlisting[language=Java, caption={frontend/src/composables/useWaypoints.js}]{../frontend/src/composables/useWaypoints.js}

\section{后端实现详解}\label{sec:backend}

\subsection{FastAPI应用架构}

\subsubsection{主应用配置}
FastAPI 应用入口 \texttt{main.py} 完成 CORS 中间件、静态文件目录与路由的装配：

\lstinputlisting[language=python, caption={backend/app/main.py}]{../backend/app/main.py}

\subsubsection{配置管理系统}
全局配置 \texttt{config.py} 管理服务地址、静态目录与 HELIOS++ 路径，所有项均支持通过环境变量覆盖：

\lstinputlisting[language=python, caption={backend/app/config.py}]{../backend/app/config.py}

\subsection{API路由实现}

\subsubsection{REST路由}
\texttt{routes.py} 实现全部 REST 端点，包含模型、航迹、配置、仿真、结果、环境诊断与缓存管理七类资源操作：

\lstinputlisting[language=python, caption={backend/app/api/routes.py}]{../backend/app/api/routes.py}

\subsection{业务服务层}

\subsubsection{航迹生成服务}
\texttt{trajectory\_generator.py} 将用户航点列表转换为HELIOS++原生.trj航迹文件，包含偏航角计算、速度驱动时间步长与转角瞬时转向逻辑：

\lstinputlisting[language=python, caption={backend/app/services/trajectory\_generator.py}]{../backend/app/services/trajectory_generator.py}

\subsubsection{点云处理服务}
\texttt{pointcloud\_parser.py} 解析LAS/LAZ与XYZ格式的点云，计算全量统计（均值/标准差/中位数/P5/P95/直方图）并提供百万级显示抽样：

\lstinputlisting[language=python, caption={backend/app/services/pointcloud\_parser.py}]{../backend/app/services/pointcloud_parser.py}

\newpage
\bibliographystyle{gbt7714-numerical}
\bibliography{cas-refs}
\newpage
\appendix

\section{部署与配置管理}

\subsection{环境变量配置}
\begin{table}[pos=!htbp]
    \centering
    \caption{环境变量配置说明}
    \label{tab:env-config}
    \begin{tabular}{@{} lll @{}}
        \toprule
        \textbf{环境变量}                  & \textbf{说明}     & \textbf{默认值}            \\
        \midrule
        \texttt{FEHALS\_HOST}          & 服务器监听地址         & 0.0.0.0                 \\
        \texttt{FEHALS\_PORT}          & 服务器端口           & 8000                    \\
        \texttt{HELIOS\_PATH}          & HELIOS++可执行文件路径 & /usr/local/bin/helios++ \\
        \texttt{HELIOS\_REPO}          & HELIOS++仓库路径    & /opt/helios++/repo      \\
        \texttt{HELIOS\_ASSETS}        & HELIOS++资源目录    & /opt/helios++/assets    \\
        \texttt{FEHALS\_SIM\_TIMEOUT}  & 仿真超时时间（秒）       & 300                     \\
        \texttt{FEHALS\_CORS\_ORIGINS} & 允许的跨域源          & http://localhost:3000   \\
        \texttt{FEHALS\_UPLOAD\_DIR}   & 文件上传目录          & uploads                 \\
        \texttt{FEHALS\_CACHE\_DIR}    & 缓存目录            & cache                   \\
        \bottomrule
    \end{tabular}
\end{table}

\subsection{Docker部署方案}
\begin{lstlisting}[language=bash,caption=Dockerfile - 容器化部署]
# Dockerfile
FROM python:3.9-slim

WORKDIR /app

# 安装系统依赖
RUN apt-get update && apt-get install -y \
    gcc \
    g++ \
    git \
    wget \
    && rm -rf /var/lib/apt/lists/*

# 安装HELIOS++
RUN wget https://github.com/3dgeo-heidelberg/helios_plusplus/releases/download/v1.0.0/helios++_linux_v1.0.0.tar.gz && \
    tar -xzf helios++_linux_v1.0.0.tar.gz && \
    mv helios++ /usr/local/bin/

# 复制requirements文件并安装Python依赖
COPY requirements.txt .
RUN pip install --no-cache-dir -r requirements.txt

# 复制应用代码
COPY . .

# 创建必要的目录
RUN mkdir -p uploads cache static logs

# 设置环境变量
ENV FEHALS_HOST=0.0.0.0
ENV FEHALS_PORT=8000
ENV HELIOS_PATH=/usr/local/bin/helios++
ENV HELIOS_REPO=/app/helios++/repo
ENV HELIOS_ASSETS=/app/helios++/assets

# 暴露端口
EXPOSE 8000

# 启动命令
CMD ["python", "run.py"]
\end{lstlisting}

\subsection{服务监控与日志}
\begin{lstlisting}[language=bash,caption=监控配置示例]
# 使用Prometheus + Grafana监控
# prometheus.yml
global:
  scrape_interval: 15s

scrape_configs:
  - job_name: 'fehals'
    static_configs:
      - targets: ['fehals-backend:8000']
    metrics_path: '/metrics'
    scrape_interval: 15s
\end{lstlisting}

\section{开发规范与最佳实践}

\subsection{代码规范}
\begin{itemize}
    \item \textbf{Python代码}：遵循PEP 8规范，使用black格式化，flake8检查
    \item \textbf{Vue代码}：使用ESLint + Prettier，采用Composition API
    \item \textbf{Git提交规范}：使用Angular提交格式（type(scope): message）
    \item \textbf{分支管理}：master（主分支）、develop（开发分支）、feature/*（功能分支）
\end{itemize}

\subsection{API设计原则}
\begin{itemize}
    \item \textbf{RESTful}：遵循REST设计规范，使用标准HTTP方法
    \item \textbf{状态码}：使用标准HTTP状态码，如200、201、400、404、500等
    \item \textbf{错误处理}：统一错误格式，包含错误代码和描述信息
    \item \textbf{文档化}：所有API接口必须有Swagger/OpenAPI文档
\end{itemize}

\section{故障排查指南}

\begin{table}[pos=!htbp]
    \centering
    \caption{常见问题及解决方案}
    \label{tab:troubleshooting}
    \begin{tabular}{@{} ll @{}}
        \toprule
        \textbf{问题现象} & \textbf{解决方案}                \\
        \midrule
        HELIOS++无法启动  & 检查HELIOS\_PATH环境变量，确认可执行文件存在 \\
        模型文件上传失败      & 检查文件格式和大小，确认在允许范围内           \\
        仿真任务卡住        & 检查HELIOS++资源目录，查看任务日志        \\
        WebSocket连接断开 & 检查CORS配置，确认代理设置正确            \\
        点云渲染性能差       & 使用降采样，减少同时渲染的点数              \\
        \bottomrule
    \end{tabular}
\end{table}

\end{document}